\section{Vorgehensweise}

Die Entwicklung von Pablo orientiert sich methodisch an der VDI-Richtlinie~2221, die ein systematisches Vorgehen zur Entwicklung technischer Produkte und Systeme beschreibt~\cite{vdi2221}.
Die Richtlinie definiert einen allgemeinen Entwicklungsprozess, der von der Klärung der Aufgabenstellung über die Erarbeitung von Funktionsstrukturen und Lösungsprinzipien bis hin zur Ausarbeitung eines Gesamtentwurfs reicht.
Ein wesentliches Merkmal dieses Vorgehensmodells ist die Möglichkeit, einzelne Arbeitsschritte iterativ zu durchlaufen.
Erkenntnisse aus späteren Phasen können so in frühere Entwurfsentscheidungen zurückfließen, ohne den gesamten Prozess neu aufsetzen zu müssen.

Dieses iterative Prinzip war für die Entwicklung von Pablo besonders geeignet, da das System mechanische, elektronische und softwaretechnische Komponenten in einem Produkt vereint.
Entwurfsentscheidungen in einem Bereich hatten häufig unmittelbare Auswirkungen auf andere Bereiche, sodass ein rein sequenzielles Vorgehen nicht praktikabel gewesen wäre.

Das Gesamtprojekt wurde daher in mehrere Entwicklungsbereiche unterteilt, an denen parallel gearbeitet wurde: Softwarearchitektur, Hardwareintegration, Gehäusedesign und Interaktionskonzept.
Diese Aufteilung ermöglichte es, die domänenspezifischen Teilaufgaben zunächst unabhängig voneinander zu bearbeiten und anschließend schrittweise zu einem funktionsfähigen Gesamtsystem zusammenzuführen.

In den folgenden Kapiteln werden die einzelnen Bereiche im Detail beschrieben.
Dabei wird jeweils das methodische Vorgehen innerhalb des Bereichs dargestellt, die erzielten Ergebnisse dokumentiert und die getroffenen Entscheidungen eingeordnet.